\section*{第 \ref{ch:g9:algebra} 章——代数与方程}

\begin{solution}{exo:g9:algebra:1}
$4(3x - 2) = 12x - 8$.

$(x+6)(x+2) = x^2 + 2x + 6x + 12 = x^2 + 8x + 12$.

$(2x-3)(x+5) = 2x^2 + 10x - 3x - 15 = 2x^2 + 7x - 15$.

$(x-1)(x-9) = x^2 - 9x - x + 9 = x^2 - 10x + 9$.
\end{solution}

\begin{solution}{exo:g9:algebra:2}
$(x+8)^2 = x^2 + 16x + 64$.

$(3x-4)^2 = 9x^2 - 24x + 16$.

$(x+10)(x-10) = x^2 - 100$.

$(5x+2)(5x-2) = 25x^2 - 4$.
\end{solution}

\begin{solution}{exo:g9:algebra:3}
$101^2 = (100+1)^2 = 10000 + 200 + 1 = 10201$.

$59^2 = (60-1)^2 = 3600 - 120 + 1 = 3481$.

$53 \times 47 = (50+3)(50-3) = 2500 - 9 = 2491$.
\end{solution}

\begin{solution}{exo:g9:algebra:4}
$9x + 12 = 3(3x + 4)$.

$x^2 - 64 = (x+8)(x-8)$.

$x^2 + 14x + 49 = (x+7)^2$.

$18x^2 - 6x = 6x(3x - 1)$.
\end{solution}

\begin{solution}{exo:g9:algebra:5}
$4x + 7 = 19$：$4x = 12$，故 $x = 3$。

$6x - 5 = 2x + 11$：$4x = 16$，故 $x = 4$。

$5(x-3) = 3x + 1$：$5x - 15 = 3x + 1$，于是 $2x = 16$ 且 $x = 8$。
\end{solution}

\begin{solution}{exo:g9:algebra:6}
$(x+4)(3x-12) = 0$：$x = -4$ 或 $x = 4$。

$x^2 - 121 = (x-11)(x+11) = 0$：$x = 11$ 或 $x = -11$。

$(2x-1)^2 = 0$：唯一解为 $x = \frac12$。

$x^2 = 8x$：$x(x - 8) = 0$，故 $x = 0$ 或 $x = 8$。
\end{solution}

\begin{solution}{exo:g9:algebra:7}
利用平方差公式，取 $a = 3x + 2$，$b = 5$：
\[
E = (3x + 2)^2 - 5^2 = (3x + 2 + 5)(3x + 2 - 5) = (3x + 7)(3x - 3).
\]
零乘积法则：由 $3x + 7 = 0$ 得 $x = -\frac73$，又由 $3x - 3 = 0$
得 $x = 1$。
\end{solution}

\begin{solution}{exo:g9:algebra:8}
$3x - 4 \geq 8$：$3x \geq 12$，故 $x \geq 4$（$4$ 处为实心点，
向右涂色）。

$5 - 2x > 1$：$-2x > -4$，两边除以 $-2$ 时不等号反向：$x < 2$
（$2$ 处为空心点，向左涂色）。

$4(x+1) \leq x - 5$：$4x + 4 \leq x - 5$，故 $3x \leq -9$，
即 $x \leq -3$（$-3$ 处为实心点，向左涂色）。
\end{solution}

\begin{solution}{exo:g9:algebra:9}
设所求数为 $x$：$3x + 7 = 5x - 9$。于是 $16 = 2x$，故 $x = 8$。
验证：$3 \times 8 + 7 = 31$，$5 \times 8 - 9 = 31$。
\end{solution}

\begin{solution}{exo:g9:algebra:10}
\emph{1.} $(n+1)(n-1) = n^2 - 1$（第三个乘法公式）。$n$ 的相邻两数
为 $n - 1$ 与 $n + 1$，故其积为 $n^2 - 1$：例如\
$24 \times 26 = 625 - 1 = 624$。

\emph{2.} 三个连续整数可记为 $n - 1$，$n$，
$n + 1$。它们的和为
\[
(n-1) + n + (n+1) = 3n ,
\]
是 $3$ 的\omterm{def:g9:arith:divisor}{倍数}（即中间那个数的三倍）。
\end{solution}

\begin{solution}{pb:g9:algebra:1}
\textbf{1.} $(10a + b) - (10b + a) = 9a - 9b = 9(a - b)$：两位数与其数字对调后的数之差
总是 $9$ 的\omterm{def:g9:arith:divisor}{倍数}。
验证：$72 - 27 = 45 = 9 \times 5$，且
$a - b = 7 - 2 = 5$。

\textbf{2.} \omterm{def:g9:algebra:expand}{展开}右边：
$9(11a + b) + a + b + c = 99a + 9b + a + b + c
= 100a + 10b + c$。因此该数等于 $9$ 的\omterm{def:g9:arith:divisor}{倍数}加上
其数字和：两者除以 $9$ 所得\omterm{def:g3:division:remainder}{余数}相同。
特别地，一者\omterm{def:g3:division:remainder}{余数}为 $0$，意味着另一者
\omterm{def:g3:division:remainder}{余数}亦为 $0$。这就是弃九法，已证。

\textbf{3.} $9(11a + b)$ 也是 $3$ 的\omterm{def:g9:arith:divisor}{倍数}，因此同一恒等式
说明该数与其数字和除以 $3$ 所得\omterm{def:g3:division:remainder}{余数}相同，
这即弃三法。

\textbf{4.} 写 $47 = 9k + 2$，$86 = 9l + 5$。于是
\[
47 \times 86 = (9k + 2)(9l + 5)
= 9\,(9kl + 5k + 2l) + 10 :
\]
乘积的\omterm{def:g3:division:remainder}{余数}与 $2 \times 5 = 10$ 相同，
即 $1$。因此正确结果的\omterm{def:g3:division:remainder}{余数}也必须为
$1$——$4\,042$ 满足。但检验只看\omterm{def:g3:division:remainder}{余数}：
$4\,042$ 与 $4\,402$（数字互换）同样通过，因此
通过并不能证明任何事情——不通过才足以定罪。

\textbf{5.} \omterm{def:g9:algebra:expand}{展开}：$11(91a + 9b + c) - a + b - c + d
= 1001a + 99b + 11c - a + b - c + d
= 1000a + 100b + 10c + d$。因此该数能被 $11$ \omterm{def:g9:arith:divisor}{整除}，
当且仅当其交替和 $-a + b - c + d$ 亦能被\omterm{def:g9:arith:divisor}{整除}。对于
$2\,926$：$-2 + 9 - 2 + 6 = 11$，可被 $11$ \omterm{def:g9:arith:divisor}{整除}，故
$2\,926 = 11 \times 266$。

\textbf{6.} $742 - 247 = 495$；$495 + 594 = 1\,089$。任一合法
例子同样得到 $1\,089$。

\textbf{7.} $(100a + 10b + c) - (100c + 10b + a) = 99(a - c)$.
对于 $a - c = 2, 3, \dots, 9$：
\[
198,\ 297,\ 396,\ 495,\ 594,\ 693,\ 792,\ 891 .
\]

\textbf{8.} $100(m - 1) + 90 + (10 - m) = 100m - 100 + 90 + 10
- m = 99m$。因此 $99m$ 的三位数字依次为 $m - 1$、$9$，
再为 $10 - m$（验证 $495$：$m = 5$ 时得 $4$、$9$、$5$）。

\textbf{9.} $99m$ 的反转数字依次为 $10 - m$、$9$，
$m - 1$，故其值为 $100(10 - m) + 90 + (m - 1)$。相加得：
\[
100\,(m - 1 + 10 - m) + 180 + (10 - m + m - 1)
= 100 \times 9 + 180 + 9 = 1\,089 ,
\]
与 $m$ 无关。无论起始数为何，魔术师
总是安全的。

\textbf{10.} 若 $a = c$，则相减得 $0$，这一戏法
立即失效。若 $a - c = 1$，则得 $99$，此时必须
视为三位字符串 $099$（反转为 $990$，
$99 + 990 = 1\,089$：得救！）。注意：差为
$1$ 时，仅当观众把相减结果
写成三位数（包括首位的零）才成立——要求差距至少为
$2$，便可避开此论证。

\textbf{11.} $2(5m + 7) + d - 14 = 10m + 14 + d - 14 =
10m + d$：十位（十月至十二月时还含百位）表示
月份，个位（以及十位）表示日期。

\textbf{12.} $2(5m + 9) + d = 10m + 18 + d$：减去 $18$。

\textbf{13.} 由第 2 题，$n - (\text{数字和 } n) =
9(11a + b)$，对三位数 $n$ 总是 $9$ 的\omterm{def:g9:arith:divisor}{倍数}
（对任意位数，类似恒等式同样成立）。$9$ 的\omterm{def:g9:arith:divisor}{倍数}的数字和
仍为 $9$ 的\omterm{def:g9:arith:divisor}{倍数}（再用第 2 题），因此观众读出的数字
必须补足到下一个
$9$ 的\omterm{def:g9:arith:divisor}{倍数}：缺失数字即为其补数。歧义在于：若所读数字之和
已是 $9$ 的\omterm{def:g9:arith:divisor}{倍数}，则划去的数字
可能为 $0$ 或 $9$——排除 $0$ 即可消除疑问。

\textbf{14.} $(10t + 5)^2 = 100t^2 + 100t + 25 =
100\,t(t + 1) + 25$：写出 $t(t+1)$，再于其后接上 $25$。
$45^2$：$4 \times 5 = 20$，故为 $2\,025$。$85^2$：
$8 \times 9 = 72$，故为 $7\,225$。$105^2$：$10 \times 11 = 110$，
故为 $11\,025$。

\textbf{15.} 五个连续整数：$n + (n+1) + (n+2) + (n+3) +
(n+4) = 5n + 10 = 5(n + 2)$：恒为 $5$ 的\omterm{def:g9:arith:divisor}{倍数}。四个
连续整数：$n + (n+1) + (n+2) + (n+3) = 4n + 6 =
4(n + 1) + 2$：其\omterm{def:g3:division:remainder}{余数}为 $2$（除数为 $4$），永不为
\omterm{def:g9:arith:divisor}{倍数}。代数能将“看来可行”化为
\emph{恒成立}，将“未找到例子”化为\emph{永不成立}——两个
任何检验都无法触及的词
（见\cref{met:g7:literal:test}）。
\end{solution}
